% SHARD-ID: RC-01-CONVENTIONS
% SHARD-TITLE: Conventions and Notation
% SHARD-SUMMARY: Fixes every notational choice used in the book, with the generated notation table from db/notation.tsv as the single source of truth.
% SHARD-SUMMARY: Records the symbol clashes found across the notes on 2026-09-12 and how each was resolved.
% SHARD-KEYWORDS: conventions, notation, normalisation, clashes, dedup, macros
% SHARD-NOTES: none
% SHARD-SCRIPTS: none

\section{Conventions and Notation}
\label{sec:conventions}

Every symbol below is a macro generated from \texttt{db/notation.tsv}; shards
may not define macros of their own, and the gate warns whenever a shard
writes a reserved control sequence (such as the bare Greek letters)
bare instead of through its macro. The table is regenerated by
\texttt{make regen} and the gate fails if it is stale.

\subsection{Standing conventions}

\begin{enumerate}
  \item \textbf{Channel normalisation.} A channel $\chan$ is normalised,
  $\chan(\dm) = \frac1{\Dk}\sum_i \Un{i}\dm\Un{i}^\dagger$, so that its
  leading eigenvalue is $1$. The unnormalised sum is $\Sig = \sum_i\Eop{i} =
  \Dk\chan$. Every Ihara--Bass statement is written with $\Sig$ or
  $\Dk\chan$; the notes that wrote $\chan$ without the $1/\Dk$ are read with
  this convention.
  \item \textbf{The size parameter $\qs$.} One letter, three roles that never
  co-occur: the field size in \cref{sec:artin-schreier}, the LPS prime in
  \cref{sec:weil-lps}, and $\Dk - 1$ for graphs and channels. In every role the
  critical modulus is $\sqrt{\qs}$. For a channel with $\Dk$ Kraus operators,
  $\qs := \Dk - 1$ is a standing identification.
  \item \textbf{Ramanujan units.} Spectral bounds are reported in adjacency
  units: a channel eigenvalue is multiplied by $\Dk = \qs + 1$ and compared
  with $2\sqrt{\qs}$. Hastings' bound $\lamH = 2\sqrt{\Dk-1}/\Dk$ is the same
  statement in channel units.
  \item \textbf{The zeta variable is $\uz$}, in $\zetaX(\uz)$,
  $\zetaT(\uz)$, $\curvezeta(\uz)$ and $\det(1-\uz\Hash)$. The ring-length
  variable of the length flow is $\ringlen$, never $\uz = t/2$.
  \item \textbf{Zeros.} A nontrivial zero is $\zero = \sigs + i\gamma$ with
  $\sigs$ its real part; $\invtemp$ is reserved for the inverse temperature of
  the Bost--Connes system. The Riemann hypothesis reads $\sigs = \frac12$.
  \item \textbf{Indices.} $n$ is the matrix size in $\Mn$ and, in the ring
  pictures, the number of sites, which for Artin--Schreier curves is the
  extension degree of $\Fqn{n}$. Walk, word and prime-cycle lengths are
  $\wl$. Counts are $\Ncount{\wl}$.
  \item \textbf{Column stacking.} $\Ad(A) = \bar A\otimes A$ as an
  $n^2\times n^2$ matrix, so $\Tr\Ad(A) = |\Tr A|^2$.
  \item \textbf{Non-backtracking operator.} $\Hash(v\otimes|i\rangle) =
  \sum_{j\ne\bari{i}}\Eop{i}v\otimes|j\rangle$: the label $i$ names the
  superoperator applied, and the next label may not be its reversal. Word
  products are composed right to left, $\Eop{i_\wl}\cdots\Eop{i_1}$, with the
  cyclic condition $i_{k+1}\ne\bari{i_k}$ for $k\in\mathbb Z/\wl$.
  \item \textbf{Artin--Schreier sign.} $\alphaw{i} = -\lambda_i(\Tmat)$, so
  that $\expsum{n} = -\sum_i\alphaw{i}^n$ for all $n$ and
  $L = \prod_i(1-\alphaw{i}\uz)$, the standard normalisation of the
  $L$-function of an exponential sum. The curve count is
  $\Ncount{n} = \qs^n + 1 - \sum_i\alphaw{i}^n$.
  \item \textbf{Canonical form.} An MPS is in (left) canonical form when
  $\sum_k\Kr{k}^\dagger\Kr{k} = 1$; then the transfer matrix has leading
  eigenvalue $1$. The spectral normalisation used in the translation note
  (leading eigenvalue $1$, rest inside the unit disc) is a consequence, not
  the definition.
  \item \textbf{Status vocabulary.} The book uses only the tags of
  \cref{sec:definitions} (\cref{def:status-vocabulary}). The
  two-value ``established / not established'' of the notes and the prior-art
  verdicts (\texttt{SAME}, \texttt{RELATED}, \texttt{NOT}) are read as
  \texttt{sketched}/\texttt{numerical} versus \texttt{open}, and as a
  provenance vocabulary confined to \cref{sec:prior-art}, respectively.
  \item \textbf{Trace.} $\Tr$ is the operator trace; the field trace is
  $\ftr$.
\end{enumerate}

\subsection{Clashes found and resolved}

The extraction of 2026-09-12 (\texttt{notes/extract/notation-and-definitions.md})
found $146$ symbols across the seven notes, with these collisions. Each
resolution below is what the table enforces.
\begin{itemize}
  \item $\psi$ carried five meanings: Chebyshev function, additive character, % bare-ok
  Lévy exponent, digamma, MPS ring state. Now: $\cheb$ is Chebyshev;
  $\addch$ the additive character; $\levy{\sigs}$ the Lévy exponent;
  $\digam$ the digamma; $\mps{n}$ the ring state.
  \item $T$ carried five: the non-backtracking operator, the length-flow
  generator, the curve-zeta variable, the Heisenberg--Weyl operator, the
  Hecke operator. Now: $\Hash$ is the non-backtracking operator; $\flowgen$
  the generator; $\uz$ the variable; $\Hei$ the Heisenberg--Weyl operator
  (only with two arguments); $\Hecke{\qs}$ the Hecke adjacency operator.
  \item $\rho$ carried three: zero, density matrix, twist. Now: $\zero$ is a % bare-ok
  zero; $\dm$ a channel input; $\hol$ a twisting representation.
  \item $\invtemp$ carried three: real part of a zero, inverse temperature,
  Deligne weight. Now: $\sigs$, $\invtemp$, $\wt$.
  \item $\cxi$ was both the completed zeta and the correlation length, within
  two pages. Now: $\cxi$ is the completed zeta; $\corrlen{k}$ the correlation
  length; $\phase{k}$ the phase.
  \item $S$ was the scattering symbol, the exponential sum and the start map.
  Now: $\Ssym$, $\expsum{n}$, $\Smap$.
  \item $J$ was the Artin--Schreier range, the reversal operator and a
  splitting matrix. Now: $\asrange$ (italic), $\Jrev$ (sans-serif),
  $\Jsplit$ (bold).
  \item $D$ was the Kraus count, the dilation unitary, a diagonal degree
  matrix and the degree operator. Now: $\Dk$, $\Dil$, and $\Dop(\uz)$; the
  diagonal degree matrix of a general graph is never needed on a bouquet.
  \item $E$ was the MPS transfer matrix and the Artin--Schreier transfer
  matrix (the same concept, kept as $\Tmat$) and the superoperators of a
  Kraus family (now $\Eop{i}$, a different glyph).
  \item $W$, $V$: the edge and vertex spaces of the non-backtracking
  construction are $\Wsp$, $\Vsp$; $\Weil(g)$ and $\Wpm$ are the Weil
  representation and its blocks; $|V|$ is a vertex count.
  \item $A$: the adjacency matrix $\Aadj$; Kraus operators $\Kr{i}$; the LPS
  Hecke adjacency $\Hecke{\qs}$.
  \item $F$: the Fourier transform $\Four$; Frobenius $\Frob$; the numerical
  window sum $\Fwin$; Weil's zero-side function $\Fweil$.
  \item $M$: a generic matrix $M$; the chirp $\chirp$.
  \item $\sigma$: the real part of $s$; the reversal involution is now % bare-ok
  $\rev$. $\Sigma$ is only ever $\Sig = \Dk\chan$. % bare-ok
  \item $\alpha$: Frobenius eigenvalues carry a subscript, $\alphaw{i}$; a % bare-ok
  quaternion is bare $\quat$.
  \item $m$: the number of Kraus pairs; walk length is $\wl$. $N$: $\dim\Vsp$;
  counts are $\Ncount{\wl}$.
  \item $\edgeev$: an edge eigenvalue; the Bost--Connes isometries are
  $\bciso{n}$. $G$: a group; the window width is $\Gwin$. $\theta$: the % bare-ok
  normal-basis generator $\nbgen$; phases are $\phase{k}$. $s$: the complex
  variable; a square root of $\qs$ mod $\pp$ is $\sqrtq$.
\end{itemize}

\subsection{Notation table}

% GENERATED by scripts/labbook_check.py --regen from db/notation.tsv. Do not edit.
{\small\begin{longtable}{p{2.3cm}p{7.6cm}p{2.4cm}p{1.3cm}}\toprule
symbol & meaning & domain & where \\ \midrule\endhead
$\operatorname{Tr}$ & trace & general & RC-01 \\
$\operatorname{Ad}$ & adjoint action, Ad(A): X -$>$ A X A\^\{\}dagger & channels & RC-01 \\
$\operatorname{Re}$ & real part & general & RC-01 \\
$\operatorname{Im}$ & imaginary part & general & RC-01 \\
$M_n$ & n x n complex matrices with the Hilbert-Schmidt inner product & channels & RC-01 \\
$\mathcal V$ & vertex space of the non-backtracking construction (V = M\_n in the channel case), N = dim V & quantum Ihara & RC-08 \\
$\mathcal W$ & edge space V (x) C\^\{\}D & quantum Ihara & RC-08 \\
$T$ & non-backtracking (Hashimoto) edge operator, classical on directed edges or quantum on the edge space & graphs, quantum Ihara & RC-02 \\
$\mathcal G$ & generator of the length flow whose trace is the prime measure, eigenvalues 1, rho, -2, -4, ... & Riemann channel & RC-03 \\
$\mathcal E_{\cdot}$ & i-th superoperator of a family (E\_i = Ad(A\_i) in the Kraus case) & quantum Ihara & RC-08 \\
$A_{\cdot}$ & i-th Kraus operator & channels & RC-01 \\
$U_{\cdot}$ & i-th unitary Kraus operator & channels & RC-01 \\
$\Sigma$ & unnormalised transfer matrix sum\_i E\_i (= D Phi for a unital channel) & channels, MPS & RC-01 \\
$\Phi$ & unital channel (1/D) sum\_i U\_i X U\_i\^\{\}dagger, or a general channel & channels & RC-01 \\
$D$ & number of Kraus operators (= degree of the bouquet); D = 2m & channels & RC-01 \\
$m$ & number of Kraus pairs, D = 2m & quantum Ihara & RC-08 \\
$\iota$ & fixed-point-free reversal involution on the index set [D], bar i = iota(i) & quantum Ihara & RC-08 \\
$\bar{\cdot}$ & reversed index iota(i) & quantum Ihara & RC-08 \\
$\mathsf J$ & backtracking operator J(v (x) $|$i$>$) = E\_i v (x) $|$bar i$>$ & quantum Ihara & RC-08 \\
$\mathsf R$ & propagate-along-edge map R: W -$>$ V & quantum Ihara & RC-08 \\
$\mathsf S$ & start-any-edge map S: V -$>$ W & quantum Ihara & RC-08 \\
$\mathcal A(u)$ & deformed adjacency u sum\_i E\_i (1 - u\^\{\}2 E\_\{bar i\} E\_i)\^\{\}\{-1\} & quantum Ihara & RC-08 \\
$\mathcal D(u)$ & deformed degree u\^\{\}2 sum\_i E\_\{bar i\} E\_i (1 - u\^\{\}2 E\_\{bar i\} E\_i)\^\{\}\{-1\} & quantum Ihara & RC-08 \\
$\zeta_{T}$ & zeta function det(1 - uT)\^\{\}\{-1\} of a non-backtracking operator & quantum Ihara & RC-08 \\
$\zeta_{\Phi}$ & quantum Ihara zeta of a unital channel with unitary Kraus operators & quantum Ihara & RC-02 \\
$\zeta_{X}$ & Ihara zeta of a graph X & graphs & RC-02 \\
$A$ & adjacency matrix of a graph & graphs & RC-02 \\
$\mathsf A_{\cdot}$ & adjacency operator of the LPS Cayley graph for norm-q quaternions (Hecke operator) & Weil-LPS & RC-05 \\
$q$ & size parameter: field size (Artin-Schreier), LPS prime (Weil-LPS), and D - 1 = degree - 1 for graphs and channels & all & RC-01 \\
$p$ & a rational prime; the modulus of the Weil representation in RC-05 & all & RC-01 \\
$\lambda_2$ & second largest eigenvalue (or singular value) of a channel or graph, normalised to leading eigenvalue 1 & expanders & RC-02 \\
$\lambda_H$ & Hastings bound 2 sqrt(D-1)/D & expanders & RC-02 \\
$\mu$ & eigenvalue of the edge operator T; pairs mu mu' = q with an adjacency eigenvalue lambda & graphs & RC-02 \\
$\ell$ & length of a walk, word or prime cycle & graphs, quantum Ihara & RC-02 \\
$N_{\cdot}$ & closed non-backtracking walk count of length ell, or point count over F\_\{q\^\{\}ell\} & graphs, curves & RC-02 \\
$P$ & prime cycle of a graph (primitive closed non-backtracking walk up to rotation) & graphs & RC-02 \\
$\mathrm{Fr}$ & Frobenius & curves & RC-06 \\
$\varrho$ & representation twisting a zeta (holonomy, Artin-Ihara twist, adjoint representation) & graphs, Deligne & RC-07 \\
$X$ & density matrix or general operator input of a channel & channels & RC-01 \\
$\rho$ & nontrivial zero of zeta, rho = sigma + i gamma & number theory & RC-01 \\
$\sigma$ & real part of s (or of a zero) & number theory & RC-01 \\
$\gamma_{\cdot}$ & imaginary part of the n-th zero & number theory & RC-01 \\
$\xi$ & completed zeta xi(s) = (1/2) s (s-1) pi\^\{\}\{-s/2\} Gamma(s/2) zeta(s) & number theory & RC-04 \\
$S$ & Lax-Phillips scattering symbol S(tau) = xi(1 - 2 i tau)/xi(1 + 2 i tau) & Riemann channel & RC-04 \\
$\tau$ & Mellin variable on the dilation group; dilation by e\^\{\}t acts as e\^\{\}\{i t tau\} & Riemann channel & RC-04 \\
$K_S$ & model space H\^\{\}2 minus S H\^\{\}2 & Riemann channel & RC-04 \\
$Z(t)$ & compressed dilation semigroup P\_\{K\_S\} e\^\{\}\{i t tau\} restricted to K\_S & Riemann channel & RC-04 \\
$\Lambda$ & von Mangoldt function & number theory & RC-01 \\
$\psi$ & Chebyshev function psi(x) = sum\_\{n $<$= x\} Lambda(n) & number theory & RC-03 \\
$\mathrm e_q$ & additive character of F\_q composed with the trace & curves & RC-06 \\
$\kappa_{\cdot}$ & Bost-Connes Levy exponent kappa\_sigma(tau) = log zeta(sigma - i tau) - log zeta(sigma) & Bost-Connes & RC-04 \\
$\nu_{\cdot}$ & von Mangoldt jump measure nu\_sigma = sum p\^\{\}\{-k sigma\}/k delta\_\{k log p\} & Bost-Connes & RC-04 \\
$\beta$ & inverse temperature of the Bost-Connes system & Bost-Connes & RC-04 \\
$\mathsf L$ & ring length variable of the length flow (prime powers sit at L = log n) & Riemann channel & RC-03 \\
$F$ & windowed zero sum F(u) of the ring-norm test & Riemann channel & RC-04 \\
$E$ & transfer matrix of a matrix product state (the Artin-Schreier transfer matrix is the instance of RC-06) & MPS, curves & RC-02 \\
$\Psi_{\cdot}$ & periodic matrix product state on n sites & MPS & RC-02 \\
$\langle\Psi_{\cdot}|\Psi_{\cdot}\rangle$ & ring norm of the periodic n-site state, = Tr E\^\{\}n & MPS & RC-02 \\
$\ell^{\mathrm{c}}_{\cdot}$ & k-th correlation length of a transfer matrix & MPS & RC-03 \\
$h$ & entropy per site (log of the leading eigenvalue) & MPS & RC-03 \\
$\mathbb F_q$ & finite field with q elements & curves & RC-01 \\
$\mathbb F_p$ & finite field with p elements & curves, Weil-LPS & RC-01 \\
$\mathbb F_{q^{\cdot}}$ & extension of degree n & curves & RC-06 \\
$\mathcal S_{\cdot}$ & exponential sum S\_n(g) = sum\_\{x in F\_\{q\^\{\}n\}\} e\_q(Tr g(x)) & curves & RC-06 \\
$J$ & range of the quadratic Artin-Schreier polynomial, dim E = q\^\{\}J & curves & RC-06 \\
$\alpha_{\cdot}$ & Frobenius eigenvalues (roots of the L-polynomial), $|$alpha$|$ = sqrt q & curves & RC-06 \\
$\mathrm{SL}_2$ & special linear group & Weil-LPS & RC-05 \\
$\mathrm{PSL}_2$ & projective special linear group & Weil-LPS & RC-05 \\
$W$ & Weil representation W(g) of SL\_2(F\_p) on l\^\{\}2(F\_p) & Weil-LPS & RC-05 \\
$W_{\pm}$ & even and odd blocks of the Weil representation, dimensions (p +- 1)/2 & Weil-LPS & RC-05 \\
$\Phi_{\cdot}$ & Weil-LPS channel for the LPS prime q (with superscript pm for the block) & Weil-LPS & RC-05 \\
$\alpha$ & integer quaternion of norm q in the LPS normalisation & Weil-LPS & RC-05 \\
$g_{\alpha}$ & image of the quaternion alpha in SL\_2(F\_p) & Weil-LPS & RC-05 \\
$Z$ & clock operator on l\^\{\}2(F\_p) & Weil-LPS & RC-05 \\
$\mathsf X$ & shift operator on l\^\{\}2(F\_p) & Weil-LPS & RC-05 \\
$\mathcal F$ & Fourier transform on l\^\{\}2(F\_p) & Weil-LPS & RC-05 \\
$\mathsf M$ & chirp operator diag e\_p(x\^\{\}2/2) & Weil-LPS & RC-05 \\
$\mathsf D$ & dilation unitary on l\^\{\}2(F\_p) & Weil-LPS & RC-05 \\
$\mathbf I$ & splitting matrix I = ((x, y), (y, -x)) with x\^\{\}2 + y\^\{\}2 + 1 = 0 mod p & Weil-LPS & RC-05 \\
$\mathbf J$ & splitting matrix J = ((0, 1), (-1, 0)) & Weil-LPS & RC-05 \\
$T(u,v)$ & Heisenberg-Weyl operator X\^\{\}u Z\^\{\}v & Weil-LPS & RC-05 \\
$c$ & pointwise weight scalar of a twisted coefficient system ($|$eigenvalue$|$ = c\^\{\}ell) & Deligne & RC-07 \\
$L$ & L-function (of a curve with coefficients, or Artin-Ihara) & graphs, curves & RC-02 \\
$(\Gamma'/\Gamma)$ & digamma function in the archimedean term of the explicit formula & number theory & RC-04 \\
$w$ & weight of a pure coefficient system ($|$Frobenius eigenvalue$|$ = q\^\{\}\{w/2\}) & Deligne & RC-07 \\
$\operatorname{tr}_{q^n/q}$ & field trace from F\_\{q\^\{\}n\} to F\_q & curves & RC-06 \\
$Z_C$ & zeta function of a curve C over F\_q, exp sum\_n N\_n u\^\{\}n / n & curves & RC-06 \\
$G_{\mathrm w}$ & Gaussian window width of the ring-norm test & Riemann channel & RC-04 \\
$F_{\mathrm W}$ & zero-side function sum\_rho e\^\{\}\{(rho - 1/2) u\} whose positive-definiteness is Weil's criterion & number theory & RC-03 \\
$\vartheta_{\cdot}$ & k-th oscillation frequency of a transfer matrix, lambda\_k = exp(-1/ell\^\{\}c\_k + i vartheta\_k) & MPS & RC-03 \\
$\theta$ & generator of a normal basis of F\_\{q\^\{\}n\} & curves & RC-06 \\
$s_q$ & a square root of q modulo p & Weil-LPS & RC-05 \\
$\mathsf m_{\cdot}$ & Bost-Connes isometry mu\_n & Bost-Connes & RC-04 \\
$\mathsf U_{\cdot}$ & Shor's unitary x -$>$ a x mod N, the level-N Galois symmetry of Bost-Connes & Bost-Connes & RC-04 \\
$\hat g$ & test window in the explicit formula (Fourier transform g) & number theory & RC-04 \\
$t$ & semigroup time in Z(t); dilation by e\^\{\}t & Riemann channel & RC-04 \\
$u$ & the zeta variable, in zeta\_X(u), zeta\_T(u), Z\_C(u), det(1 - uT) & all & RC-01 \\
$\ell^2(\cdot)$ & Hilbert space of square-summable functions on a set & general & RC-01 \\
$\mathcal O_{\mathrm{Dw}}$ & Dwork's p-adic transfer operator, multiplication by the splitting function then decimation & curves & RC-06 \\
$\operatorname{dec}_{\qs}$ & decimation x\^\{\}m -$>$ x\^\{\}\{m/q\} on p-adic power series & curves & RC-06 \\
$\ell\text{-adic}$ & ell-adic (cohomology); this ell is not a walk length & Deligne & RC-06 \\
$\mathcal X$ & generator of the geodesic flow on Gamma\textbackslash\{\}G, X = H/2 through left-invariant fields; (e\^\{\}\{-tX\} f)(m) = f(phi\_\{-t\} m) & Selberg & RC-09B \\
$\mathsf H$ & left-invariant field of H = diag(1,-1) & Selberg & RC-09B \\
$\mathsf E$ & left-invariant field of E = ((0,1),(1,0)) & Selberg & RC-09B \\
$\mathsf W$ & left-invariant field of W = ((0,1),(-1,0)), the compact direction & Selberg & RC-09B \\
$\Omega$ & Casimir (1/4)(H\^\{\}2 + E\^\{\}2 - W\^\{\}2) & Selberg & RC-09B \\
$\Delta$ & positive hyperbolic Laplacian -y\^\{\}2(d\_x\^\{\}2 + d\_y\^\{\}2) & Selberg & RC-09B \\
$\lambda_{\cdot}$ & j-th Laplace eigenvalue, lambda\_j = 1/4 + r\_j\^\{\}2 & Selberg & RC-09B \\
$r_{\cdot}$ & spectral parameter of lambda\_j; r\_0 = i/2 for the constant mode, r\_j = i nu\_j for exceptional eigenvalues & Selberg & RC-09B \\
$Z_{\mathrm{Sel}}$ & Selberg zeta prod\_gamma prod\_\{k$>$=0\} (1 - e\^\{\}\{-(s+k) ell\_gamma\}) & Selberg & RC-09B \\
$\zeta_{R}$ & Ruelle zeta prod\_gamma (1 - e\^\{\}\{-s ell\_gamma\}) = Z\_Sel(s)/Z\_Sel(s+1) & Selberg & RC-09B \\
$P_{\gamma}$ & linearised Poincare map of the closed geodesic gamma & Selberg & RC-09B \\
$\operatorname{Tr}^{\flat}$ & Guillemin flat trace & Selberg & RC-09B \\
$\varsigma$ & Laplace variable of the flat trace (the resonance variable), never the Selberg variable s & Selberg & RC-09B \\
$\varLambda_{\mathrm{fl}}$ & Laplace transform of the flat trace, int\_0\^\{\}infty e\^\{\}\{-varsigma t\} Tr\^\{\}flat e\^\{\}\{-tX\} dt & Selberg & RC-09B \\
$D_{\mathrm{tow}}$ & the tower prod\_\{j$>$=1\} Z\_Sel(varsigma + j), the flat dynamical determinant & Selberg & RC-09B \\
$\mathcal B_0$ & nonconstant first band \{-1/2 +- i r\_j : j $>$= 1\} & Selberg & RC-09B \\
$\varphi$ & scattering determinant of PSL(2,Z), sqrt(pi) Gamma(s-1/2) zeta(2s-1)/(Gamma(s) zeta(2s)) & modular & RC-09B \\
$m_{\mathrm{sc}}$ & scattering multiplier -(1/2)(varphi'/varphi)(1/2 + i r) & modular & RC-09B \\
$C_{\mathrm{cusp}}$ & inverse Fourier transform of the scattering multiplier & modular & RC-09B \\
$P_{+}$ & positive-time prime measure sum\_n Lambda(n)/n delta(t - 2 log n) & modular & RC-09B \\
$P_{\mathrm{ch}}$ & prime measure sum\_n Lambda(n) n\^\{\}\{-1/2\} delta(t - 2 log n) of the Riemann-channel note & modular & RC-09B \\
$A_{\mathrm{arch}}$ & archimedean-plus-boundary distribution 1/2 - 1/(4 sinh($|$t$|$/2)) away from zero & modular & RC-09B \\
\bottomrule\end{longtable}}

